This paper presents CMWSL, a public-data Causal Multiscale Wavelet Spillover Learning framework for stock-index volatility forecasting and risk early warning. The framework was designed to answer a narrow but important question: when does a causal multiscale representation add useful predictive information beyond strong persistence benchmarks? The answer emerging from the evidence is clear and disciplined.

First, daily index volatility remains strongly persistence-driven. HAR is not a weak comparator but the correct baseline, and it remains the strongest average model in the main Rogers--Satchell specification. Second, within-index multiscale wavelet representation becomes more valuable as the forecasting problem becomes less myopic. The gains are more visible at $h=5$ and $h=10$, in richer public-information environments, and under the Parkinson target. Third, the spillover extension provides genuine incremental information for broad-market indices. Clark--West tests detect significant gains in five of nine index--horizon cells, with the strongest evidence for the S\&P~500 at $h=10$ (CW\,$=4.83$, $p<0.001$). Tail-conditioned and rolling-window diagnostics further show that these gains are concentrated in upper-volatility regimes and synchronized stress windows rather than in average conditions.

For financial risk management practice, these results carry a clear operational message. Under normal market conditions, HAR-class persistence models remain the most reliable and interpretable baseline for daily volatility risk measurement. The CMWSL extension is most valuable as a medium-horizon risk-monitoring overlay---activated by elevated macro-financial stress signals, widening credit spreads, or synchronized cross-index volatility clustering---where its frequency-specific spillover and multiscale features capture risk dynamics that a purely linear model misses. The integrated early-warning classifier provides a directly auditable, threshold-calibrated alert mechanism that connects the forecasting pipeline to actionable risk management decisions.

Several extensions are natural. International index panels, cross-asset systems, longer robustness windows, and formal model-confidence-set analysis would all strengthen external validity. More specialized sequence models may also be worth testing if they are evaluated under the same chronology and reproducibility constraints. For the present paper, however, the main conclusion is already well supported: causal multiscale wavelet spillover learning offers a useful, transparent, and reproducible extension to persistence-based market risk monitoring, and its value is most visible---and most consequential for risk management practice---when cross-market stress becomes more synchronized, more state-dependent, and more challenging for linear risk models to track in real time.
